\section{第\ref{sec:serocomputer}章　\nt{SERO}ニューロン}

\nt{SERO}ニューロン自体の性質（リズム、自己抑制、受容体による符号、サブシステム）は直接測定されている。本章の計算（調節器、フィルタ、論理）は本章の方程式から導かれる。$\tau_c$、セロトニンの拡散係数、放出部位の数は推定である。脳でループがレベル調節器として働くことは仮説である（縫線核では自己抑制は一対一で、短い休止を生むだけである）。\nt{SERO}が時間幅を担うことも仮説だが、これを支持する行動実験がある。

{\footnotesize
\setlength{\tabcolsep}{3pt}\renewcommand{\arraystretch}{1.15}
\begin{longtable}{>{\raggedright}p{0.5cm}>{\raggedright}p{4.0cm}>{\raggedright}p{5.6cm}>{\raggedright}p{2.3cm}>{\raggedright\arraybackslash}p{1.7cm}}
\toprule
No. & 主張 & 実験と測定内容 & 出典 & 状態 \\
\midrule\endfirsthead
\toprule
No. & 主張 & 実験と測定内容 & 出典 & 状態 \\
\midrule\endhead
1 & セロトニンの作用の符号は受け手の受容体が選ぶ（命題~\ref{prop:receptor}） & セロトニン受容体は薬理と機構によってファミリーに分けられる。5-HT$_{1A}$はGタンパク質を介してカリウムチャネルを開き、細胞を抑制する。5-HT$_{2A}$とイオンチャネル型の5-HT$_3$は興奮させる。一つの伝達物質が細胞ごとに逆の応答を生む。 & \cite{BarnesSharp1999} & 測定済み \\
2 & 覚醒時の\nt{SERO}ニューロンは毎秒数回のスパイクを一定に出す & 自由行動下のネコの背側縫線核ニューロンの記録。静かな覚醒時に$2.8\pm0.2$スパイク/sのゆっくりした規則的発火。徐波睡眠で発火率は$34$--$68\,\%$下がり、レム睡眠では$98\,\%$下がる。麻酔下のラットでは$1.7\pm0.5$スパイク/sで、非常に規則的。 & \cite{TrulsonJacobs1979, GagnonParent2014, JacobsAzmitia1992} & 測定済み \\
3 & \nt{SERO}ニューロンはペースメーカーである。スパイクごとの後押しは要らないが、一定の背景入力は要る（スライスでは$\alpha_1$受容体経由のノルアドレナリン、生体では\nt{NORA}から） & 脳スライスでは細胞はゆっくり規則的に発火し、各スパイクの後に$200$--$800\ms$の後過分極がある。スライスで自発発火する細胞は生体より少ない。規則的なリズムには$\alpha_1$受容体を介したノルアドレナリンの持続的入力が必要で、その遮断でリズムは消える。 & \cite{VandermaelenAghajanian1983} & 測定済み \\
4 & 受容体はほぼすべて代謝型。応答は遅く、およそ1秒のうちに立ち上がって減衰する & 5-HT$_3$を除くすべての受容体ファミリーはGタンパク質共役型である。スライスでは、誘発されたセロトニン放出が5-HT$_{1A}$を介して抑制電流を生み、それは1秒以内に立ち上がって減衰する。 & \cite{BarnesSharp1999, CourtneyFord2016} & 測定済み \\
5 & 投射先領域ではセロトニンは間隙だけでなく体積に出る。縫線核自体では近隣の抑制は一対一 & スライスでの高速サイクリックボルタンメトリー。1スパイクあたりの放出量は、シナプスのある領域と、シナプスがほとんどない縫線核とで同じであり、取り込みや受容体の遮断に依存しない。終末あたりの分子数はトランスポーターや受容体の数よりはるかに少なく、細胞外濃度は主要受容体の親和性と一致する。縫線核では5-HT$_{1A}$を介した抑制は一対一で、トランスポーターがセロトニンを間隙外に溜めない。 & \cite{Bunin1998, CourtneyFord2016} & 測定済み \\
6 & \eqref{eq:seroneuron}の濃度は回収によって減る。$\tau_c$が1秒程度というのは推定 & 生体脳でのボルタンメトリー。細胞外セロトニンを主に制限するのは合成ではなく、再取り込みと分解である。引用した研究に$\tau_c$の直接測定はなく、その大きさは本章の推定である。 & \cite{Hashemi2012serotonin} & 測定済み。$\tau_c$は推定 \\
7 & ペースメーカー1個でNOTとNOR。\nt{SERO}論理の完全性（命題~\ref{prop:serocomplete}、図~\ref{fig:serogates}） & \eqref{eq:seroneuron}から導かれる。抑制されるペースメーカーはNORであり、NORは関数的に完全である。\nt{SERO}ニューロンで論理を組んだ実験はない。体積が分かれているという条件は脳では満たされない。 & 本章で導出 & 理論 \\
8 & セロトニンは\nt{SERO}ニューロン自身をその上の受容体を介して抑制する & 麻酔下ラットで、5-HT$_{1A}$作動薬は静脈内投与でもイオントフォレシスでも、セロトニン自体と同じ用量反応関係で背側縫線核ニューロンの発火を抑制する。5-HT$_{1B}$作動薬はほとんど効かない。スライスでは、近隣細胞の内因性セロトニンが5-HT$_{1A}$を介して電流と発火の短い休止を生むが、平均発火率は変わらない。 & \cite{Sprouse1987, CourtneyFord2016} & 測定済み \\
9 & モデルでは、共通の体積を通るループはレベル調節器であり、ばらつきと時定数を$1+G$分の1にする~\eqref{eq:feedback}。脳でループがそう働くことは仮説 & 負帰還増幅器の古典的な式で、本章で導き直した。\nt{SERO}系のループゲイン$G$は測定されていない。測定されているのは、スライスでの自己抑制が一対一で、短い休止を生み、平均発火率を変えないことである。 & \cite{Black1934feedback, CourtneyFord2016} & 理論。脳では仮説 \\
10 & 濃度は発火率の移動平均、ローパスフィルタ~\eqref{eq:lowpass}、図~\ref{fig:lowpass} & 線形方程式~\eqref{eq:seroneuron}の解。図は$\tau_c=1\,\text{s}$でこの式から計算した。\texttt{models/}に個別のモデルはない。 & 本章で導出 & 理論 \\
11 & 間隙の屈曲は小分子の拡散をおよそ$2.6$倍遅くする & 点源法。スライスと生体脳でテトラメチルアンモニウムをイオントフォレシスで与え、その濃度をイオン選択性電極で記録する。屈曲度$\lambda\approx1.6$、つまり実効$D$は自由拡散の$\lambda^2\approx2.6$分の1。細胞外体積の割合は約$0.2$。セロトニン自体の組織内拡散係数はこれらの研究では測定されておらず、本章では推定である。 & \cite{Sykova2008} & 測定済み \\
12 & 独立な「配線」の長さ$\ell\sim\sqrt{D\tau}\approx20$~\textmu m~\eqref{eq:diffusionlength}。配線の数はそのような領域の数で制限される（命題~\ref{prop:volumewires}） & 拡散の法則に$D$と$\tau$の推定値（6行と11行）を代入したもの。命題~\ref{prop:volumewires}はその帰結である。容積伝達の独立チャネルを数えた直接の実験はない。 & 本章で導出 & 理論 \\
13 & 1個の\nt{SERO}ニューロンの出力は脳の広大な範囲に広がる。放出部位は推定$\sim10^5$ & ラット背側縫線核の単一ニューロンの完全再構成。上行軸索はすべて強く分岐し、1細胞の軸索の総長は最大$18.7$~cm、1細胞の枝が線条体、前頭前皮質、扁桃体に届く。細胞あたりの放出部位数は直接数えられていない。 & \cite{GagnonParent2014} & 測定済み。$10^5$は推定 \\
14 & \nt{SERO}ニューロンは出力と応答の異なるいくつかのサブシステムに分かれる & マウスでのウイルス遺伝学的標識。扁桃体に投射するニューロンと前頭皮質に投射するニューロンは分岐がほとんど重ならず、異なる入力を受け、不快刺激に逆向きに応答する。各群の活性化と抑制は行動を異なる形で変える。 & \cite{Ren2018} & 測定済み \\
15 & 指数割引での忍耐の条件$D<\tau_\gamma\ln(R/r_0)$~\eqref{eq:patience}。測定される双曲割引では$D<\tau_\gamma(R/r_0-1)$ & 割引$e^{-D/\tau_\gamma}$または$1/(1+D/\tau_\gamma)$から導かれる。導出は本章。秒単位の遅延での選択課題のサル。割引は指数より双曲線に近い。遅延報酬の予告信号に対する\nt{DOPA}ニューロンの応答も、遅延とともにやはり双曲線的に下がる。 & 本章で導出、\cite{KobayashiSchultz2008} & 理論 \\
16 & \nt{SERO}が時間幅$\tau_\gamma$を設定する（\ref{sec:horizon}節） & マウスで背側縫線核の\nt{SERO}ニューロンを光遺伝学的に活性化すると、遅延報酬を待つ時間が延び、まれになった報酬の探索の粘り強さが増す。ヒトではセロトニンの低下（トリプトファン欠乏食）で遅延報酬の割引が急になる。濃度がどのように$\tau_\gamma$に変わるかは分かっていない。 & \cite{Doya2002, Miyazaki2014, Lottem2018, Schweighofer2008} & 仮説 \\
\bottomrule
\end{longtable}}
